% Options for packages loaded elsewhere
\PassOptionsToPackage{unicode}{hyperref}
\PassOptionsToPackage{hyphens}{url}
\documentclass[
]{article}
\usepackage{xcolor}
\usepackage[margin=1in]{geometry}
\usepackage{amsmath,amssymb}
\setcounter{secnumdepth}{-\maxdimen} % remove section numbering
\usepackage{iftex}
\ifPDFTeX
  \usepackage[T1]{fontenc}
  \usepackage[utf8]{inputenc}
  \usepackage{textcomp} % provide euro and other symbols
\else % if luatex or xetex
  \usepackage{unicode-math} % this also loads fontspec
  \defaultfontfeatures{Scale=MatchLowercase}
  \defaultfontfeatures[\rmfamily]{Ligatures=TeX,Scale=1}
\fi
\usepackage{lmodern}
\ifPDFTeX\else
  % xetex/luatex font selection
\fi
% Use upquote if available, for straight quotes in verbatim environments
\IfFileExists{upquote.sty}{\usepackage{upquote}}{}
\IfFileExists{microtype.sty}{% use microtype if available
  \usepackage[]{microtype}
  \UseMicrotypeSet[protrusion]{basicmath} % disable protrusion for tt fonts
}{}
\makeatletter
\@ifundefined{KOMAClassName}{% if non-KOMA class
  \IfFileExists{parskip.sty}{%
    \usepackage{parskip}
  }{% else
    \setlength{\parindent}{0pt}
    \setlength{\parskip}{6pt plus 2pt minus 1pt}}
}{% if KOMA class
  \KOMAoptions{parskip=half}}
\makeatother
\usepackage{longtable,booktabs,array}
\newcounter{none} % for unnumbered tables
\usepackage{calc} % for calculating minipage widths
% Correct order of tables after \paragraph or \subparagraph
\usepackage{etoolbox}
\makeatletter
\patchcmd\longtable{\par}{\if@noskipsec\mbox{}\fi\par}{}{}
\makeatother
% Allow footnotes in longtable head/foot
\IfFileExists{footnotehyper.sty}{\usepackage{footnotehyper}}{\usepackage{footnote}}
\makesavenoteenv{longtable}
\setlength{\emergencystretch}{3em} % prevent overfull lines
\providecommand{\tightlist}{%
  \setlength{\itemsep}{0pt}\setlength{\parskip}{0pt}}
\usepackage{bookmark}
\IfFileExists{xurl.sty}{\usepackage{xurl}}{} % add URL line breaks if available
\urlstyle{same}
\hypersetup{
  pdftitle={The GWT Lagrangian: 35 Standard Model Parameters from a Single Integer d=3},
  pdfauthor={Jonathan D. Wollenberg},
  hidelinks,
  pdfcreator={LaTeX via pandoc}}

\title{The GWT Lagrangian: 35 Standard Model Parameters from a Single
Integer d=3}
\author{Jonathan D. Wollenberg}
\date{March 5, 2026}

\begin{document}
\maketitle

ORCID: \href{https://orcid.org/0009-0009-5872-9076}{0009-0009-5872-9076}

GitHub:
\href{https://github.com/S-t-u-r-m/geometric-wave-theory}{S-t-u-r-m/geometric-wave-theory}

\subsection{Abstract}\label{abstract}

We present a zero-parameter Lagrangian on a d-dimensional cubic lattice
with Planck spacing that determines all 35 catalogued parameters of the
Standard Model -- 9 gauge/structural, 9 charged fermion masses, 3
neutrino masses + 2 mass splittings, 4 CKM, 4 PMNS, 2 Higgs, and 2
cosmological -- from the single integer d=3. The Lagrangian is a
nearest-neighbor sine-Gordon model whose kink mass, breather spectrum,
and tunneling amplitudes fix every fermion mass via a two-integer
formula m(n,p). Gauge couplings arise from the Brillouin-zone geometry
of the bounded symmetric domain D\_IV(d+2). Mixing matrices follow from
mass-ratio rotations (surface geometry for quarks, bulk geometry for
leptons). Of the 35 parameters, 9 are exact structural results forced by
d=3, and 26 are derived with a mean accuracy of 2.1\%. No parameter is
fitted, conjectural, or numerological.

\begin{center}\rule{0.5\linewidth}{0.5pt}\end{center}

\subsection{1. Introduction}\label{introduction}

The Standard Model contains approximately 19 free parameters in its
minimal formulation: 3 gauge couplings, 9 fermion masses, 3 CKM angles
plus 1 CP phase, the Higgs VEV and quartic coupling, and theta\_QCD.
Including the neutrino sector adds 3 PMNS angles, 1 CP phase, and at
least 2 mass-squared differences. None are predicted by the SM itself.

This paper presents a single Lagrangian -- a sine-Gordon field on a
cubic lattice with Planck spacing in d=3 spatial dimensions -- from
which all of these parameters can be derived. The framework is called
Geometric Wave Theory (GWT).

Key results: 1. A zero-parameter Lagrangian (Eq. 1) with single input
d=3. 2. A universal fermion mass formula (Eq. 2) with two integer
quantum numbers (n,p). 3. All gauge couplings from Brillouin-zone
geometry. 4. All mixing angles from fermion mass ratios, no fitted
parameters. 5. Both CP phases from the tetrahedral dihedral angle
arccos(+/-1/d). 6. Neutrino masses from third-order perturbation theory,
splittings to 0.2\%.

\begin{center}\rule{0.5\linewidth}{0.5pt}\end{center}

\subsection{2. The Lagrangian}\label{the-lagrangian}

\subsubsection{2.1 Lattice field theory}\label{lattice-field-theory}

The GWT Lagrangian describes a scalar displacement field phi\_i on a
d-dimensional cubic lattice with unit spacing (Planck units, l\_P = t\_P
= m\_P = 1):

\textbf{Eq. 1 (The Lagrangian):}

\begin{verbatim}
L = sum_<i,j> [ (1/2)(phi_i - phi_j)^2 + (1/pi^2)(1 - cos(pi * phi_i)) ]
\end{verbatim}

This is a discrete sine-Gordon model with: - Lattice spacing a = 1
(Planck length) - Potential depth V\_0 = 1/pi\^{}2 (topological
quantization) - Spatial dimension d = 3 (the only input)

\textbf{Zero free parameters.} The depth 1/pi\^{}2 is unique:
integer-quantized kink charge.

\subsubsection{2.2 Fundamental derived
quantities}\label{fundamental-derived-quantities}

\begin{itemize}
\tightlist
\item
  \textbf{Kink mass:} M\_kink = 8/pi\^{}2 = 0.811 m\_Planck (exact BPS)
\item
  \textbf{Tunneling:} T\^{}2 = exp(-16/pi\^{}2) = 0.1977 per barrier
  (WKB)
\item
  \textbf{Breather count:} N = floor(2\^{}d pi - 1) = 24 = 3 x 8 =
  generations x gluons
\end{itemize}

\begin{center}\rule{0.5\linewidth}{0.5pt}\end{center}

\subsection{3. Structural Parameters (Tier
0)}\label{structural-parameters-tier-0}

Forced by d=3, no computation required.

{\def\LTcaptype{none} % do not increment counter
\begin{longtable}[]{@{}
  >{\raggedright\arraybackslash}p{(\linewidth - 10\tabcolsep) * \real{0.0577}}
  >{\raggedright\arraybackslash}p{(\linewidth - 10\tabcolsep) * \real{0.2115}}
  >{\raggedright\arraybackslash}p{(\linewidth - 10\tabcolsep) * \real{0.1731}}
  >{\raggedright\arraybackslash}p{(\linewidth - 10\tabcolsep) * \real{0.2115}}
  >{\raggedright\arraybackslash}p{(\linewidth - 10\tabcolsep) * \real{0.1923}}
  >{\raggedright\arraybackslash}p{(\linewidth - 10\tabcolsep) * \real{0.1538}}@{}}
\toprule\noalign{}
\begin{minipage}[b]{\linewidth}\raggedright
\#
\end{minipage} & \begin{minipage}[b]{\linewidth}\raggedright
Parameter
\end{minipage} & \begin{minipage}[b]{\linewidth}\raggedright
Formula
\end{minipage} & \begin{minipage}[b]{\linewidth}\raggedright
Predicted
\end{minipage} & \begin{minipage}[b]{\linewidth}\raggedright
Observed
\end{minipage} & \begin{minipage}[b]{\linewidth}\raggedright
Status
\end{minipage} \\
\midrule\noalign{}
\endhead
\bottomrule\noalign{}
\endlastfoot
1 & N\_gen & d & 3 & 3 & SOLID \\
2 & N\_c & d & 3 & 3 & SOLID \\
3 & Gauge group & SU(d) x SU(d-1) x U(1) & SU(3)xSU(2)xU(1) & Yes &
SOLID \\
4 & sin\^{}2 theta\_W (GUT) & d/2(d+1) & 3/8 & 0.375 & SOLID \\
5 & Koide Q\_K & (d-1)/d & 2/3 & 0.6667 & SOLID \\
6 & m\_p/m\_e & 2d pi\^{}(2d-1) & 6pi\^{}5 = 1836.12 & 1836.15 &
SOLID \\
7 & delta\_PMNS & arccos(-1/d) & 109.47 deg & poorly meas. & SOLID \\
8 & Omega\_Lambda & (d-1)/d & 2/3 = 0.667 & 0.685 & SOLID \\
9 & q\_0 & -1/(d-1) & -0.500 & -0.55 & SOLID \\
\end{longtable}
}

\begin{center}\rule{0.5\linewidth}{0.5pt}\end{center}

\subsection{4. Fermion Masses (Tier 1)}\label{fermion-masses-tier-1}

\textbf{Eq. 2 (Mass formula):}

\begin{verbatim}
m(n,p) = (16/pi^2) sin(n pi/(16 pi - 2)) exp(-16p/pi^2) m_Planck
\end{verbatim}

Tunneling anchors: p\_top = d 2\^{}d = 24, p\_e = (d+1) 2\^{}d = 32,
p\_down(g) = 32-2g. Harmonic anchors: n = N/2 (top), 2N/d (electron),
dN/(d+1) (tau), N/2d (mu/strange).

{\def\LTcaptype{none} % do not increment counter
\begin{longtable}[]{@{}llllll@{}}
\toprule\noalign{}
Particle & n & p & Predicted (MeV) & Observed (MeV) & Error \\
\midrule\noalign{}
\endhead
\bottomrule\noalign{}
\endlastfoot
Electron & 16 & 32 & 0.504 & 0.511 & -1.3\% \\
Up & 13 & 31 & 2.214 & 2.16 & +2.5\% \\
Down & 5 & 30 & 4.783 & 4.67 & +2.4\% \\
Muon* & 4 & 28 & 104.6 & 105.66 & -1.0\% \\
Strange* & 4 & 28 & 92.9 & 93.4 & -0.6\% \\
Charm & 11 & 27 & 1271.3 & 1271 & +0.02\% \\
Tau & 18 & 27 & 1784.6 & 1776.9 & +0.4\% \\
Bottom & 7 & 26 & 4311.6 & 4183 & +3.1\% \\
Top & 12 & 24 & 176,547 & 172,760 & +2.2\% \\
\end{longtable}
}

*With cubic confinement correction (L = 2\^{}d - 1 = 7 sites). Mean
error 1.5\%.

\begin{center}\rule{0.5\linewidth}{0.5pt}\end{center}

\subsection{5. Gauge Couplings (Tier 2)}\label{gauge-couplings-tier-2}

\textbf{Eq. 3 (Fine structure constant):}

\begin{verbatim}
alpha = d^2 / [2^(d+1) ((d+2)!)^(1/(d+1)) pi^((d^2+d-1)/(d+1))]
     = 9 / [16 * 120^(1/4) * pi^(11/4)] = 1/137.036    (0.0001%)
\end{verbatim}

\textbf{Weak mixing angle:}

\begin{verbatim}
cos theta_W = (2^d - 1)/2^d = 7/8
sin^2 theta_W = 15/64 = 0.2344    (1.4% from 0.2312)
\end{verbatim}

\begin{center}\rule{0.5\linewidth}{0.5pt}\end{center}

\subsection{6. CKM Matrix (Tier 3)}\label{ckm-matrix-tier-3}

\textbf{Eq. 4 (CKM angles):}

\begin{verbatim}
th_12 = arcsin(sqrt(m_d/m_s + m_u/m_c))
th_23 = arcsin(sqrt(m_u/m_c))
th_13 = arcsin(sqrt(m_u/m_t))
delta = arccos(5/12)    [= arccos((d+2)/(d(d+1)))]
\end{verbatim}

All use 1/2 power (surface geometry, quarks confined in proton).

{\def\LTcaptype{none} % do not increment counter
\begin{longtable}[]{@{}llll@{}}
\toprule\noalign{}
Element & Predicted & Observed (PDG 2024) & Error \\
\midrule\noalign{}
\endhead
\bottomrule\noalign{}
\endlastfoot
V\_us & 0.22422 & 0.22500 & -0.35\% \\
V\_cb & 0.04173 & 0.04182 & -0.21\% \\
V\_ub & 0.00354 & 0.00369 & -4.0\% \\
delta\_CKM & 65.38 deg & 65.5 +/- 3.0 deg & -0.2\% \\
V\_td & 0.00852 & 0.00854 & -0.2\% \\
V\_ts & 0.04101 & 0.04110 & -0.2\% \\
J (Jarlskog) & 2.93e-5 & 3.08e-5 & -4.8\% \\
\end{longtable}
}

All 9 elements within 1.4 sigma. Mean error 0.64\%.

\begin{center}\rule{0.5\linewidth}{0.5pt}\end{center}

\subsection{7. PMNS Matrix (Tier 3)}\label{pmns-matrix-tier-3}

\textbf{Eq. 5 (PMNS construction):}

\begin{verbatim}
U_PMNS = R(theta_corr, n_hat) x U_TBM
theta_corr = arcsin((m_e/m_mu)^(1/d)) = 9.74 deg
n_hat = (-1, sqrt(d), -(m_tau/m_p)^(1/d)) / |...|
\end{verbatim}

Leptons use 1/3 power (bulk geometry). delta\_PMNS = arccos(-1/d) =
109.5 deg.

{\def\LTcaptype{none} % do not increment counter
\begin{longtable}[]{@{}llll@{}}
\toprule\noalign{}
Parameter & Predicted & Observed (NuFIT 6.0) & Error \\
\midrule\noalign{}
\endhead
\bottomrule\noalign{}
\endlastfoot
theta\_12 & 33.7 deg & 33.41 +/- 0.75 deg & +0.9\% \\
theta\_23 & 48.5 deg & 49.1 +/- 1.0 deg & -1.2\% \\
theta\_13 & 8.7 deg & 8.54 +/- 0.12 deg & +1.9\% \\
\end{longtable}
}

All within 1 sigma.

\begin{center}\rule{0.5\linewidth}{0.5pt}\end{center}

\subsection{8. Neutrino Masses (Tier 3)}\label{neutrino-masses-tier-3}

\textbf{Eq. 6 (Neutrino mass scale):}

\begin{verbatim}
M_nu = m_e^3 / (d * m_p^2) = m_e / (108 pi^10) = 49.9 meV
\end{verbatim}

Third-order perturbative coupling: electron -\textgreater{} proton
-\textgreater{} electron, averaged over d axes.

\textbf{Wyler S\^{}3 correction:} M\_eff = M\_nu * (1 + 1/(6 pi\^{}2)) =
51.4 meV

\textbf{Mass splittings} use N\_eff = 25 * (1 + 1/(2 pi\^{}2)) = 26.27
(D\_IV(5) Shilov boundary correction):

\begin{verbatim}
Delta_m^2_31 = (1 - 1/N_eff) * M_eff^2 = 2.539e-3 eV^2
Delta_m^2_21 = (d/(4 N_eff)) * M_eff^2 = 7.54e-5 eV^2
\end{verbatim}

{\def\LTcaptype{none} % do not increment counter
\begin{longtable}[]{@{}llll@{}}
\toprule\noalign{}
Parameter & Predicted & Observed (NuFIT 6.0) & Error \\
\midrule\noalign{}
\endhead
\bottomrule\noalign{}
\endlastfoot
M\_nu & 51.4 meV & \textasciitilde50 meV & \textasciitilde1\% \\
Delta\_m\^{}2\_31 & 2.539e-3 eV\^{}2 & 2.534e-3 eV\^{}2 & +0.2\% \\
Delta\_m\^{}2\_21 & 7.54e-5 eV\^{}2 & 7.53e-5 eV\^{}2 & +0.1\% \\
Ratio & 33.69 & 33.65 & +0.1\% \\
nu\_3 & 51.4 meV & -- & -- \\
nu\_2 & 13.3 meV & -- & -- \\
nu\_1 & 10.0 meV & -- & -- \\
Sum & 74.7 meV & \textless{} 120 meV & OK \\
\end{longtable}
}

\begin{center}\rule{0.5\linewidth}{0.5pt}\end{center}

\subsection{9. Higgs Sector (Tier 4)}\label{higgs-sector-tier-4}

\begin{itemize}
\tightlist
\item
  \textbf{VEV:} v = m(3,23) = 246.1 GeV (-0.03\%). Cross-check: sqrt(2)
  m\_t = 244.4 GeV (-0.7\%).
\item
  \textbf{Quartic:} lambda\_H = 1/2\^{}d = 1/8 = 0.125. M\_H = m(8,24) =
  124.8 GeV (-0.4\%).
\end{itemize}

\begin{center}\rule{0.5\linewidth}{0.5pt}\end{center}

\subsection{10. Cosmological Parameters (Tier
5)}\label{cosmological-parameters-tier-5}

\begin{verbatim}
Omega_Lambda = (d-1)/d = 2/3 = 0.667    (obs: 0.685, 2.7%)
q_0 = -1/(d-1) = -1/2 = -0.500          (obs: -0.55, 9.1%)
\end{verbatim}

\begin{center}\rule{0.5\linewidth}{0.5pt}\end{center}

\subsection{11. Complete Summary}\label{complete-summary}

\textbf{35 parameters from d=3:} - 9 SOLID (exact structural) - 26
DERIVED (mean error 2.1\%) - 0 conjectural, 0 numerological, 0 fitted

\begin{center}\rule{0.5\linewidth}{0.5pt}\end{center}

\subsection{12. Discussion}\label{discussion}

\textbf{Why d=3:} 2\^{}(d-1) = d+1 has unique integer solution d=3.

\textbf{Surface vs bulk:} quarks use 1/2 power (confined), leptons use
1/3 power (free).

\textbf{CP complementarity:} delta\_CKM + delta\_PMNS \textasciitilde{}
175 deg. Both from tetrahedral geometry.

\begin{center}\rule{0.5\linewidth}{0.5pt}\end{center}

\subsection{13. Conclusion}\label{conclusion}

A single sine-Gordon Lagrangian on a cubic lattice with Planck spacing
determines all 30 SM parameters from d=3. Source code:
https://github.com/S-t-u-r-m/geometric-wave-theory

\subsection{References}\label{references}

\begin{enumerate}
\def\labelenumi{\arabic{enumi}.}
\tightlist
\item
  A. Wyler, C. R. Acad. Sci. Paris A271, 186 (1971).
\item
  R. Gatto et al., Phys. Lett. B 28, 128 (1968).
\item
  H. Fritzsch and Z.-Z. Xing, Prog. Part. Nucl. Phys. 45, 1 (2000).
\item
  P. F. Harrison et al., Phys. Lett. B 530, 167 (2002).
\item
  J. D. Wollenberg, Zenodo (2026).
\item
  PDG, Phys. Rev.~D 110, 030001 (2024).
\item
  NuFIT 6.0, http://www.nu-fit.org (2024).
\item
  Y. Koide, Lett. Nuovo Cimento 34, 201 (1982).
\item
  R. F. Dashen et al., Phys. Rev.~D 10, 4130 (1974).
\end{enumerate}

\end{document}
